\cvsection{Expériences professionnelles}

\begin{cventries}
  \cventry
    {Stagiaire ingénieur de recherche en IA}
    {INRIA}
    {Bordeaux, France}
    {Mars 2026 -- Septembre 2026}
    {
      \begin{cvitems}
        \item {Analyse d’un besoin de recherche d’articles scientifiques et conception en autonomie d’une application intégrant un agent IA multi-étapes.}
        \item {Construction de l’agent sans framework d’orchestration avec API OpenAI, conception de prompts et gestion des outils, de l’état, de la mémoire et du contexte.}
        \item {Mise en œuvre d’un pipeline RAG avec chunking, embeddings, recherche sémantique et base vectorielle ChromaDB.}
        \item {Développement d’un backend Python/FastAPI avec API REST documentées avec OpenAPI et persistance des documents et métadonnées dans MongoDB.}
        \item {Développement d’une interface React/Next.js permettant aux chercheurs d’interagir avec l’agent, ses outils et de sélectionner les messages conservés dans le contexte.}
        \item {Automatisation de la collecte de métadonnées scientifiques et transmission des résultats de tâches cron par webhooks.}
        \item {Gestion d’une application web conteneurisée permettant d’utiliser des LLM installés localement avec LM Studio sur un Mac Studio, avec installation et configuration des modèles.}
        \item {Utilisation de vLLM pour lancer et servir des modèles sur une infrastructure HPC équipée de GPU NVIDIA.}
        \item {Évaluation de Qwen, Gemma et GLM et de plusieurs quantifications avec Slurm selon la latence, le TTFT, la qualité et la taille du contexte.}
        \item {Définition du protocole expérimental, analyse des résultats et rédaction en cours d’un article scientifique comparant les modèles.}
        \item {Conteneurisation avec Docker, gestion des sources avec Git et mise en place d’une CI/CD sur GitHub.}
      \end{cvitems}
    }

  \cventry
    {Stagiaire ingénieur de recherche en apprentissage par renforcement}
    {Université Karunya}
    {Coimbatore, Inde}
    {Juin 2025 -- Sept. 2025}
    {
      \begin{cvitems}
        \item {Développement d’un modèle d’apprentissage par renforcement pour la navigation autonome d’un robot hospitalier, puis déploiement et tests sur Raspberry Pi.}
      \end{cvitems}
    }
\end{cventries}
